% !TeX root = ../applied-mathematics1-companion.tex
% Dimensionless illustration N=1, S0=.99, I0=.01, rho=.3; not observed data.
\begin{center}
\begin{tikzpicture}[x=7.4cm,y=5cm,>=Stealth]
\draw[->] (0,0)--(1.06,0) node[right] {$S/N$};
\draw[->] (0,0)--(0,1.06) node[above] {$I/N$};
\draw[muted,dotted] (0,1)--(1,0);
\draw[muted,dashed] (.3,0)--(.3,.7);
\draw[very thick,ink,domain=0.0404097352:.99,samples=100]
 plot (\x,{1-\x+.3*ln(\x/.99)});
\draw[->,thick,ink] (.72,{1-.72+.3*ln(.72/.99)})
 --(.66,{1-.66+.3*ln(.66/.99)});
\draw[->,thick,ink] (.15,{1-.15+.3*ln(.15/.99)})
 --(.12,{1-.12+.3*ln(.12/.99)});
\fill (.99,.01) circle (1.5pt) node[above left,yshift=2pt] {initial};
\fill (.3,{.7+.3*ln(.3/.99)}) circle (1.5pt) node[above right] {peak};
\node[below] at (.3,0) {$\rho/N$};
\node[above right,muted] at (.68,.32) {$S+I=N$};
\node[below left] at (0,0) {$0$};
\end{tikzpicture}

{\small Illustrative SIR trajectory; arrows move towards decreasing $S$.}
\end{center}
